
%%%%%%%%%%%%%%%%%%%%%%%%%%%%%%%%%%%%%%%%
%           main body
%%%%%%%%%%%%%%%%%%%%%%%%%%%%%%%%%%%%%%%%

%-----------------------------------
%	TITLE PAGE
%-----------------------------------

\title[第五章\\
函数空间Lp]{\texorpdfstring{\LARGE \bfseries 第五章\quad 函数空间 $L^p$}{第五章 函数空间Lp}} % The short title appears at the bottom of every slide, the full title is only on the title page
\author[\S 5.2\\
Lp空间可分性] % Your institution as it will appear on the bottom of every slide, maybe shorthand to save space
{\Large \bfseries \S 5.2 $L^p$ 空间的可分性}
% \author{Singular Value Decomposition} % Your name
% \date{\today} % Date can be changed to a custom date
\date{}

%------------------------------------------------

\begin{document}

%------------------------------------------------

\begin{frame}

\titlepage % Print the title page as the first slide

\end{frame}

%------------------------------------------------

\section[引言]{引言: 完备性之后}

%------------------------------------------------

\begin{frame}{引言: 从 $\mathbb{Q} \subset \mathbb{R}$ 说起}

\begin{itemize}
\item 上一节证明了 $L^p$ 空间的\alert{完备性}; $L^p$ 空间还有一个比较重要的性质:\alert{可分性}, 存在\alert{可列}的稠密子空间;
\item 类比: $\mathbb{Q}$ 是 $\mathbb{R}$ 中的\alert{可列稠密子集}, 有理数足够``稠'', 又足够``少'' (可列), 于是 $\mathbb{R}$ 的结构可以通过
  $\mathbb{Q}$ 来把握;
\item 有了可分性, $L^p$ 空间的结构就比较简单.
\end{itemize}

\pause

\begin{center}
\fbox{我们这一节最重要的结论就是: $L^p$ 是可分空间 ($1 \leqslant p < \infty$).}
\end{center}

\pause

\begin{itemize}
\item 路线: 主线为\alert{三步逼近引理} (有界可测 $\to$ 简单 $\to$ 有理系数);
  支线为 Weierstra\ss{} 逼近 (多项式, 超纲补充); 收尾为\alert{弱收敛}及收敛关系总结.
\end{itemize}

\end{frame}

%------------------------------------------------

\section[稠密与可分]{稠密与可分}

%------------------------------------------------

\begin{frame}{定义: 稠密与可分}

\begin{Def}[稠密]
设 $A$ 是 $L^p$ 空间的一个子集. 若 $\forall f \in L^p$, 恒有 $A$ 中的点列
$\{ g_n \}_{n \in \mathbb{N}}$, 使得 $g_n \xrightarrow{\ \text{强}\ } f$,
则称 $A$ 在 $L^p$ 中\alert{稠密}.
\end{Def}

\pause

\begin{Def}[可分]
若存在\alert{可列}子集 $A$, 使得 $A$ 在 $L^p$ 中稠密, 则称 $L^p$ 是\alert{可分的}.
\end{Def}

\pause

\begin{itemize}
\item \textbf{注}: 以上定义中的 $L^p$ 空间可以替换为一般的度量空间,
  得到一般的\alert{可分度量空间}的定义;
\item \textbf{类比}: $\mathbb{Q}$ 可列且在 $\mathbb{R}$ 中稠密
  $\Longrightarrow$ $\mathbb{R}$ (按绝对值距离) 可分.
\end{itemize}

\end{frame}

%------------------------------------------------

\section[可分定理]{$L^p$ 可分性定理}

%------------------------------------------------

\begin{frame}{定理: $L^p$ 可分 ($1 \leqslant p < \infty$)}

\begin{thm}[$L^p$ 的可分性]
设 $1 \leqslant p < \infty$, 则 $L^p$ 是可分空间.
\end{thm}

\pause

\begin{itemize}
\item \textbf{证明的思路}: 需要一些引理, 用一些\alert{特殊的函数}来逼近一般的函数;
\item \textbf{回忆}: 在学习单个可积函数性质时 (\S 4.2 附录), 我们证明过
  可积函数可用连续函数\alert{平均 ($L^1$) 逼近}, 这里是更细致一些的结论.
\end{itemize}

\pause

\begin{remark}
定理断言中要求 $p < \infty$: $L^\infty$ 不可分
(如 $\big\{ \chi_{[n, n+1)} \big\}_{n \in \mathbb{N}}$ 两两距离为 $1$).
\end{remark}

\end{frame}

%------------------------------------------------

\begin{frame}{引理: 三步逼近}

\begin{lem}[三步逼近]
设 $f \in L^p(E)$, ($1 \leqslant p < \infty$), 那么 $\forall \varepsilon > 0$:
\begin{enumerate}
\item 存在\alert{有界可测}函数 $g$, 使得 $\lVert f - g \rVert_p < \varepsilon$;
\item 设 $g$ 为\alert{有界可测集 $E$ 上的有界可测}函数, 那么存在\alert{简单函数}
  $\varphi$, 使得 $\lVert g - \varphi \rVert_p < \varepsilon$;
\item 设 $\varphi$ 为有界可测集 $E$ 上的简单函数, 则存在形如
\[
s(x) = \sum_{k=1}^{m} r_k \, \chi_{e_k}(x)
\qquad \text{(全体记为 } S\text{)}
\]
的简单函数, 其中 $r_k \in \mathbb{Q}$, $e_k$ 为互不相交的、端点为有理数的开区间,
使得 $\lVert \varphi - s \rVert_p < \varepsilon$.
\end{enumerate}
\end{lem}

\pause

\begin{itemize}
\item 三个函数类层层收窄: 一般 $\to$ 有界 $\to$ 阶梯 $\to$ \alert{有理数据}.
\end{itemize}

\end{frame}

%------------------------------------------------

\section[三步逼近]{三步逼近引理的证明}

%------------------------------------------------

\begin{frame}{(1) 的证明: 渐缩集列}

\startproof[bg=yellow, fg=red]((1) 的证明)
由 $|f|^p \in L_E$ 及积分的\alert{绝对连续性}, $\forall \varepsilon > 0$,
$\exists \delta > 0$, 使得 $\forall$ 可测集 $e \subset E$, 只要 $me < \delta$,
就有 $\int_e |f|^p ~ \mathrm{d} m < \varepsilon^p$.

\pause
考虑 $E$ 中的渐缩集列
\[
A_k = E\big( |f| > k \big),
\qquad
A_1 \supset A_2 \supset \cdots \supset A_k \supset \cdots
\]
\begin{itemize}
\item 由于 $|f|^p \in L_E$: $\forall k \in \mathbb{N}$, $mA_k < \infty$
  (特别地 $mA_1 < \infty$);
\item 又由于 $f$ 在 $E$ 上 a.e.\ 有限, 所以
\[
m \Big( \bigcap_{k=1}^{\infty} A_k \Big)
= m \big( E\big( |f| = \infty \big) \big) = 0 ;
\]
\item 从而 $\lim_k mA_k = 0$ (测度的上连续性): 对 $\delta > 0$,
  $\exists k_0 \in \mathbb{N}$, 使得 $\forall k > k_0$, 有 $mA_k < \delta$.
\end{itemize}
\qed

\end{frame}

%------------------------------------------------

\begin{frame}{(1) 的证明: 截断与估计}

\startproof[bg=yellow, fg=red]((1) 的证明, 续)
令
\[
g(x) =
\begin{cases}
f(x), & \text{若 } x \in E \setminus A_{k_0}, \\
0, & \text{若 } x \in A_{k_0}.
\end{cases}
\]
\pause
那么 $g$ 是有界可测函数 ($|g| \leqslant k_0$), 且有
\[
\lVert f - g \rVert_p
= \Big( \int_E |f - g|^p ~ \mathrm{d} m \Big)^{\!1/p}
= \Big( \int_{A_{k_0}} |f|^p ~ \mathrm{d} m \Big)^{\!1/p}
< \big( \varepsilon^p \big)^{1/p}
= \varepsilon .
\]
\qed

\pause

\begin{itemize}
\item \textbf{注意}: (1) 不要求 $mE < \infty$, 绝对连续性对一般可测集上的可积函数成立;
\item 与 \S 4.3 有界收敛定理的``截大值''手法相同: 在 $E(|f| > k_0)$ 上截断.
\end{itemize}

\end{frame}

%------------------------------------------------

\begin{frame}{(2) 的证明: 正负部与简单函数逼近}

\startproof[bg=yellow, fg=red]((2) 的证明)
设 $E$ 有界可测, $g$ 有界可测, 设一个上界为 $M$ ($|g| \leqslant M$).
写 $g = g_+ - g_-$, 那么 $g_+, g_-$ 非负、有界、可测, 且可积.

\pause
由非负可测函数积分的定义
($\int_E h ~ \mathrm{d} m = \sup \big\{ \int_E \varphi ~ \mathrm{d} m : 0 \leqslant \varphi \leqslant h, \varphi \text{ 简单} \big\}$),
$\forall \varepsilon > 0$, 存在 $E$ 上非负简单函数
\[
0 \leqslant \varphi_1 \leqslant g_+,
\qquad
0 \leqslant \varphi_2 \leqslant g_-,
\]
使得
\[
\int_E |g_+ - \varphi_1| ~ \mathrm{d} m < \frac{\varepsilon^p}{2 \cdot (2M)^{p-1}},
\qquad
\int_E |g_- - \varphi_2| ~ \mathrm{d} m < \frac{\varepsilon^p}{2 \cdot (2M)^{p-1}} .
\]
\pause
令 $\varphi = \varphi_1 - \varphi_2$, 即有 $\varphi$ 为有界简单函数
($|\varphi_1|, |\varphi_2| \leqslant M$, 从而 $|\varphi| \leqslant M$).
\qed

\end{frame}

%------------------------------------------------

\begin{frame}{(2) 的证明: $(2M)^{p-1}$ 技巧}

\startproof[bg=yellow, fg=red]((2) 的证明, 续)
首先
\[
\int_E |g - \varphi| ~ \mathrm{d} m
\leqslant \int_E \big( |g_+ - \varphi_1| + |g_- - \varphi_2| \big) ~ \mathrm{d} m
< \frac{\varepsilon^p}{(2M)^{p-1}} .
\]
由于 $|g| \leqslant M$, $|\varphi| \leqslant M$, 所以 $|g - \varphi| \leqslant 2M$, 从而
\[
\int_E |g - \varphi|^p ~ \mathrm{d} m
\leqslant (2M)^{p-1} \int_E |g - \varphi| ~ \mathrm{d} m
< (2M)^{p-1} \cdot \frac{\varepsilon^p}{(2M)^{p-1}}
= \varepsilon^p ,
\]
故 $\lVert g - \varphi \rVert_p < \varepsilon$.
\qed

\pause

\begin{itemize}
\item \textbf{技巧的核心}: $|h| \leqslant 2M$ 时 $|h|^p \leqslant (2M)^{p-1} |h|$, ``降一次幂'', 而 $\varphi_i$ 的逼近精度恰好带一个 $(2M)^{p-1}$ 分母, 两者相消;
\item 这也是选取 $\varepsilon^p / \big( 2(2M)^{p-1} \big)$ 这一``奇怪''精度的原因.
\end{itemize}

\end{frame}

%------------------------------------------------

\begin{frame}{(3) 的证明: 闭集逼近与互不相交开集}

\startproof[bg=yellow, fg=red]((3) 的证明)
设 $\varphi(x) = \sum_{i=1}^{n} c_i \chi_{E_i}(x)$,
$E_i \subset E$ 为互不相交可测集, 且 $E = \bigcup_i E_i$;
记 $M = \max_i |c_i|$.

\pause
\begin{itemize}
\item \textcircled{1} \textbf{闭集逼近} (可测集的正则性): $\forall \varepsilon > 0$,
  $\forall E_i$, 存在闭集 $F_i \subset E_i$, 使得
\[
m \big( E_i \setminus F_i \big) < \frac{\varepsilon^p}{2 n \, (2M)^p}
\qquad (i = 1, \dots, n);
\]
\item \textcircled{2} \textbf{互不相交开邻域}: $F_i$ 为互不相交闭集,
  所以存在互不相交开集 $\widetilde{G}_i \supset F_i$.
  ($n=2$ 情形: 令 $\rho = \rho(F_1, F_2) > 0$,
  取 $\widetilde{G}_1 = \bigcup_{y \in F_1} U(y, \rho/3)$,
  $\widetilde{G}_2 = \bigcup_{x \in F_2} U(x, \rho/3)$ 即可.)
\end{itemize}
\qed

\end{frame}

%------------------------------------------------

\begin{frame}{(3) 的证明: 有限化与端点有理化}

\startproof[bg=yellow, fg=red]((3) 的证明, 续)
\begin{itemize}
\item \textcircled{3} \textbf{有限化}: 又由于 $F_i$ 是有界闭集,
  $\widetilde{G}_i$ 的构成区间对 $F_i$ 构成开覆盖, 从而可选出有限覆盖,
  将其并集记为 $G_i$. 那么
  \[
  F_i \subset G_i \subset \widetilde{G}_i
  \quad \text{为开集, 且 } G_i \text{ 只有有限多个 (互不相交的) 构成区间};
  \]
\item \textcircled{4} \textbf{端点有理化}: 可以假设 $G_i$ 的构成区间的端点都是
  有理数, 否则只需将每个构成区间的端点向外微调到充分近的有理数
  (外扩幅度不超过 $\min_{i \neq j} \rho(F_i, F_j) / 4$ 量级),
  互不相交性不变.
\end{itemize}

\pause
\begin{itemize}
\item \textcircled{5} \textbf{有理系数}: $\forall c_i$, 取 $r_i \in \mathbb{Q}$, 使得
  \[
  |r_i - c_i| < \frac{\varepsilon}{(2 \, mE)^{1/p}}
  \qquad \text{且} \qquad |r_i| \leqslant |c_i|
  \quad \Longrightarrow \quad
  |\varphi - s| \leqslant 2M ;
  \]
  并令
  \[
  s(x) = \sum_{i=1}^{n} r_i \, \chi_{G_i}(x) .
  \]
\end{itemize}
\qed

\end{frame}

%------------------------------------------------

\begin{frame}{(3) 的证明: 误差估计}

\startproof[bg=yellow, fg=red]((3) 的证明, 续)
由 $G_i$ 互不相交: 在 $F_i$ 上 $\varphi = c_i$ 且 $s = r_i$ (无混叠), 于是
\[
\int_E |\varphi - s|^p ~ \mathrm{d} m
= \int_{\bigcup_i E_i} |\varphi - s|^p ~ \mathrm{d} m
= \Big( \int_{\bigcup_i F_i} + \int_{\bigcup_i (E_i \setminus F_i)} \Big)
| \varphi - s |^p ~ \mathrm{d} m
\]
\pause
\[
\leqslant \Big( \frac{\varepsilon}{(2 \, mE)^{1/p}} \Big)^{\!p}
m \Big( \bigcup_{i=1}^{n} F_i \Big)
+ (2M)^p \; m \Big( \bigcup_{i=1}^{n} (E_i \setminus F_i) \Big)
\leqslant \frac{\varepsilon^p}{2 \, mE} \cdot mE
+ (2M)^p \sum_{i=1}^{n} \frac{\varepsilon^p}{2n (2M)^p}
\]
\[
= \frac{\varepsilon^p}{2} + \frac{\varepsilon^p}{2}
= \varepsilon^p
\qquad \Longrightarrow \qquad
\lVert \varphi - s \rVert_p < \varepsilon .
\]
\qed

\pause

\begin{itemize}
\item 两项误差各占一半: $\bigcup F_i$ 上靠\alert{系数有理化},
  $\bigcup (E_i \setminus F_i)$ 上靠\alert{测度控制}与 $|\varphi - s| \leqslant 2M$.
\end{itemize}

\end{frame}

%------------------------------------------------

\section[可分性证明]{可分性定理的证明}

%------------------------------------------------

\begin{frame}{定理的证明: $\varepsilon/3$ 拼接与 $S$ 可列}

\startproof[bg=yellow, fg=red]($L^p$ 可分性定理的证明)
设 $1 \leqslant p < \infty$, 任取 $f \in L^p$, $\varepsilon > 0$. 由引理:
\[
\exists\ \text{有界可测 } g: \lVert f - g \rVert_p < \varepsilon/3;
\quad
\exists\ \text{简单函数 } \varphi: \lVert g - \varphi \rVert_p < \varepsilon/3;
\quad
\exists\ s \in S: \lVert \varphi - s \rVert_p < \varepsilon/3 .
\]
于是
\[
\lVert f - s \rVert_p
\leqslant \lVert f - g \rVert_p + \lVert g - \varphi \rVert_p
+ \lVert \varphi - s \rVert_p
< \varepsilon .
\]
\pause
而
\[
S = \Big\{ s = \sum_{i=1}^{m} r_i \chi_{[\alpha_i, \beta_i]}
:\ r_i,\ \alpha_i,\ \beta_i \in \mathbb{Q},\ m \in \mathbb{N} \Big\}
\]
是可列集 ($\mathbb{Q}^3$ 可数, 有限长序列之并可数).
故 $S$ 是 $L^p$ 中的可列稠密集, $L^p$ 可分.
\qed

\pause

\begin{itemize}
\item \textbf{注}: 构成区间的端点开闭不影响 (有理端点处至多差一个单点集, 零测).
\end{itemize}

\end{frame}

%------------------------------------------------

\section[正算子]{另一条路: 多项式逼近}

%------------------------------------------------

\begin{frame}{多项式逼近: 正算子与 Bernstein 算子}

(接下来我们用\alert{另一种方法}证明 $L^p$ 空间的可分性,
即 Weierstra\ss{} 逼近定理, 用多项式进行逼近. 先引入一些概念.)

\pause

\begin{Def}[正算子与多项式算子]
称线性算子 $L: C([a,b]) \longrightarrow C([a,b])$ 为\alert{正算子},
若 $\forall$ 非负 $f \in C([a,b])$, 总有 $L(f)$ 为非负连续函数;
称 $L$ 为\alert{多项式算子}, 若 $\forall f \in C([a,b])$, $L(f)$ 为多项式.

\vspace{0.3em}
\textbf{注}: 正算子保持不等号: \
$L(f - g) \geqslant 0 \Longrightarrow L(f) \geqslant L(g)$.
\end{Def}

\pause

\begin{eg}[Bernstein 算子]
$C([0,1])$ 上的线性算子
\[
f \longmapsto B_n(f; x) := \sum_{k=0}^{n} f \Big( \frac{k}{n} \Big)
\binom{n}{k} x^k (1 - x)^{n-k}
\]
为\alert{正的多项式算子}; $B_n(f; x)$ 被称为 $f$ 的 \alert{Bernstein 多项式}.
\end{eg}

\end{frame}

%------------------------------------------------

\section[Korovkin]{Korovkin 定理}

%------------------------------------------------

\begin{frame}{Korovkin 定理}

\begin{attnblock}[超纲补充]
以下 Korovkin 定理及其证明为超出大纲的补充内容, 供感兴趣的同学选读;
其作用是把 Bernstein 定理归结为对三个单项式的验证.
\end{attnblock}

\graymode

\pause

\begin{thm}[Korovkin]
设 $\{ L_n \}_{n \in \mathbb{N}}$ 为 $C([a,b])$ 上的一列正线性算子,
且对 $i = 0, 1, 2$ 满足 $L_n(\alpha_i; x)$ \alert{一致收敛}于 $\alpha_i$,
其中 $\alpha_i(x)$ 为单项式 $x^i$.
那么 $\forall f \in C([a,b])$, $L_n(f; x)$ 一致收敛于 $f$.
\end{thm}

\pause

\begin{itemize}
\item \textbf{证明的路线}: 一致连续性给出逐点控制不等式
  $\to$ 正算子\alert{保序}地``传递''到算子值
  $\to$ 三个单项式的假设控制每一项 $\to$ 取 $z = x$ 收尾;
\item 关键辅助函数: $\varphi(t) = (t - x)^2$, 由 $1, t, t^2$ 线性表出.
\end{itemize}

\normalmode

\end{frame}

%------------------------------------------------

\begin{frame}{Korovkin 的证明 (一): 逐点控制}

\graymode

\startproof[bg=yellow, fg=red](Korovkin 定理的证明)
由于 $f \in C([a,b])$ 为连续函数, $[a,b]$ 为紧集, 所以 $f$ 一致连续:
$\forall \varepsilon > 0$, $\exists \delta > 0$, 使得 $\forall x, t \in [a,b]$,
只要 $|x - t| < \delta$, 就有 $|f(x) - f(t)| < \varepsilon / 2$.

\pause
令 $M = \max_{t \in [a,b]} |f(t)|$ (连续函数在闭区间上能取到最值),
$\varphi(t) = (t - x)^2$. 那么 $\forall x, t \in [a,b]$, 下列不等式恒成立:
\[
|f(t) - f(x)| < \varepsilon + \frac{2M}{\delta^2} \, \varphi(t) :
\]
\begin{itemize}
\item 若 $|t - x| < \delta$: $|f(t) - f(x)| < \varepsilon/2
  < \varepsilon + \frac{2M}{\delta^2} \varphi(t)$;
\item 若 $|t - x| \geqslant \delta$: $\varphi(t) = (t-x)^2 \geqslant \delta^2$
  $\Longrightarrow \frac{2M}{\delta^2} \varphi(t) \geqslant 2M
  \geqslant |f(t) - f(x)|$.
\end{itemize}
\qed

\normalmode

\end{frame}

%------------------------------------------------

\begin{frame}{Korovkin 的证明 (二): 作用正算子}

\graymode

\startproof[bg=yellow, fg=red](Korovkin 定理的证明, 续)
将 $t$ 视为变元, $x$ 视为常值, 将 $L_n$ 作用于上式两边 (去绝对值后按两边夹).
由于 $L_n$ 为正的线性算子, 保持不等号不变, 有
\[
- \varepsilon \, L_n(1; z) - \frac{2M}{\delta^2} L_n(\varphi; z)
\leqslant L_n(f; z) - f(x) \, L_n(1; z)
\leqslant \varepsilon \, L_n(1; z) + \frac{2M}{\delta^2} L_n(\varphi; z),
\]
\pause
其中由线性性,
\[
L_n(\varphi; z) = L_n(t^2 - 2xt + x^2; z)
= L_n(t^2; z) - 2x \, L_n(t; z) + x^2 \, L_n(1; z)
\xrightarrow{\ \text{一致}\ } z^2 - 2xz + x^2 \cdot 1 = (z - x)^2 .
\]
\pause
移项 (补上 $f(x) = f(x) \cdot 1$) 可得
\[
\big| L_n(f; z) - f(x) \big|
\leqslant \varepsilon \, L_n(1; z)
+ \frac{2M}{\delta^2} L_n(\varphi; z)
+ |f(x)| \, \big| L_n(1; z) - 1 \big| .
\tag{$*$}
\]
\qed

\normalmode

\end{frame}

%------------------------------------------------

\begin{frame}{Korovkin 的证明 (三): 一致收尾}

\graymode

\startproof[bg=yellow, fg=red](Korovkin 定理的证明, 续)
由 $L_n(1; z) \rightrightarrows 1$ 与 $L_n(\varphi; z) \rightrightarrows (z-x)^2$
(关于 $x, z \in [a,b]$ 一致), 存在与 $x, z$ 无关的 $N \in \mathbb{N}$,
使得 $\forall n > N$:
\[
\big| L_n(1; z) - 1 \big| < \varepsilon
\ \Big( \text{从而 } L_n(1; z) < 1 + \varepsilon \Big),
\qquad
\big| L_n(\varphi; z) - (z - x)^2 \big| < \varepsilon ,
\ \forall z \in [a,b] \ \text{成立}.
\]
\pause
特别地, 当 $z = x$ 时 (此时 $(z-x)^2 = 0$, 故 $L_n(\varphi; x) < \varepsilon$),
\[
(*) \leqslant \varepsilon (1 + \varepsilon) + \frac{2M}{\delta^2} \cdot \varepsilon
+ M \cdot \varepsilon
\leqslant \Big( 2 + \frac{2M}{\delta^2} + M \Big) \varepsilon
\qquad (0 < \varepsilon < 1),
\]
右端与 $x$ 无关. 这说明了 $L_n(f; \cdot)$ 在 $[a,b]$ 一致收敛于 $f$.
\qed

\pause

\begin{itemize}
\item 常规的``$\varepsilon$-板''论证: 对每个 $\varepsilon$, $n > N$ 后误差被
  \alert{常数倍} $\varepsilon$ 控制, 令 $\varepsilon \to 0$ 即得一致收敛于 $0$.
\end{itemize}

\normalmode

\end{frame}

%------------------------------------------------

\section[逼近定理]{Bernstein 与 Weierstra\ss{}}

%------------------------------------------------

\begin{frame}{推论: Bernstein 定理}

\begin{cor}[Bernstein 定理]
$\forall f \in C([0,1])$, Bernstein 多项式
\[
B_n(f; x) = \sum_{k=0}^{n} f \Big( \frac{k}{n} \Big) \binom{n}{k}
x^k (1 - x)^{n-k}, \qquad n \in \mathbb{N},
\]
在 $[0,1]$ 上一致收敛于 $f$.
\end{cor}

\pause

\startproof[bg=yellow, fg=red](证明)
$B_n$ 是正线性算子, 只需对单项式 $f = 1, x, x^2$ 进行验证 (Korovkin 定理):
\[
B_n(1; x) = 1 = f(x);
\qquad
B_n(t; x) = x = f(x);
\qquad
B_n(t^2; x) = x^2 + \frac{x (1 - x)}{n} \longrightarrow x^2 = f(x),
\]
其中 $\dfrac{x(1-x)}{n} \leqslant \dfrac{1}{4n} \to 0$ 对 $x \in [0,1]$ 一致.
\qed

\end{frame}

%------------------------------------------------

\begin{frame}{推论: Weierstra\ss{} 逼近定理}

\begin{cor}[Weierstra\ss{} 逼近定理]
设 $f \in C([a,b])$, 那么 $\forall \varepsilon > 0$, 存在多项式 $P(x)$,
使得 $\forall x \in [a,b]$, 有 $|f(x) - P(x)| < \varepsilon$.
\end{cor}

\startproof[bg=yellow, fg=red](证明)
先经变量替换 $t = \dfrac{x - a}{b - a}$ 化到 $[0,1]$:
$g(t) := f\big( a + (b - a) t \big) \in C([0,1])$.

考虑 Bernstein 多项式 $B_n(g; t)$, 其一致收敛于 $g$:
即 $\forall \varepsilon > 0$, 存在与 $t$ 无关的 $N \in \mathbb{N}$,
使得 $\forall n > N$,
\[
\big| B_n(g; t) - g(t) \big| < \varepsilon
\qquad \text{对 } t \in [0,1] \text{ 恒成立}.
\]
取多项式 $P(x) := B_n \Big( g;\ \dfrac{x - a}{b - a} \Big)$ 即可.
\qed

\end{frame}

%------------------------------------------------

\begin{frame}{由此又得 $L^p$ 可分}

\begin{itemize}
\item \textbf{注}: 由 Weierstra\ss{} 逼近定理, 可以推出 $L^p$ 空间的可分性.
\end{itemize}

\[
L^p \supset \{ \text{有界可测函数} \}
\overset{\text{Luzin}}{\supset} \{ \text{连续函数} \}
\supset \{ \text{多项式} \}
\supset \{ \text{有理系数多项式} \}
\]
\vspace{-0.5em}
\begin{itemize}
\item 有界可测函数稠密: 三步逼近引理的 (1) (2);
\item 连续函数稠密: Luzin 定理 ($mE < \infty$ 时按 $L^p$ 逼近);
\item 多项式稠密: Weierstra\ss{} 逼近定理;
\item 有理系数多项式稠密: 系数有理化 ($\mathbb{Q}$ 在 $\mathbb{R}$ 中稠密),
  且它只有\alert{可列}多个, 即得可分.
\end{itemize}

\pause

\begin{remark}
此路线适用于 $E$ 为(含于)有限区间的情形; 三步逼近的路线则更一般.
两条路线殊途同归: $1 \leqslant p < \infty$ 时 $L^p$ 可分.
\end{remark}

\end{frame}

%------------------------------------------------

\section[弱收敛]{弱收敛}

%------------------------------------------------

\begin{frame}{定义: 弱收敛}

(在 $L^p$ 空间中, 除了依\alert{范数}的强收敛概念之外, 还有\alert{弱收敛}的概念:
强收敛只依赖于自己 (范数), 弱收敛还依赖于与其它函数的配对.)

\pause

\begin{Def}[弱收敛]
设 $1 \leqslant p < \infty$, $q$ 为 $p$ 的\alert{相伴数}
($\frac{1}{p} + \frac{1}{q} = 1$), $f, f_n \in L^p$, $n \in \mathbb{N}$.
若 $\forall g \in L^q$, 有
\[
\lim_{n \to \infty} \int_E f_n g ~ \mathrm{d} m = \int_E f g ~ \mathrm{d} m ,
\]
则称 $\{ f_n \}$ 在 $L^p$ 中\alert{弱收敛}于 $f$,
记为 $f_n \xrightarrow{\ \text{弱}\ } f$.
\end{Def}

\end{frame}

%------------------------------------------------

\begin{frame}{双线性配对与对偶问题}

\begin{itemize}
\item \textcircled{1} 可以定义\alert{双线性映射}
  \[
  L^p \times L^q \longrightarrow \mathbb{R},
  \qquad
  (f, g) \longmapsto \int_E f \cdot g ~ \mathrm{d} m .
  \]
  固定 $g$, 令 $\varphi_g(f) := \int_E f g ~ \mathrm{d} m$,
  则 $\varphi_g$ 是 $L^p$ 上的线性泛函:
  $\sim$ $L^q \subset (L^p)^*$, $g \mapsto \varphi_g$.
\end{itemize}

\pause

\begin{question}[两个问题]
\begin{itemize}
\item (a) $\lVert g \rVert_q \geqslant \lVert \varphi_g \rVert
  := \sup_{\lVert f \rVert_p = 1} \Big| \int_E f g ~ \mathrm{d} m \Big|$, Hölder 给出"$\leqslant$"方向; \alert{等号}成立吗?
\item (b) $L^q = (L^p)^*$ 吗? 一般\alert{不}相等
  ($L^q$ 只是弱收敛``看得见''的部分).
\end{itemize}
\end{question}

\end{frame}

%------------------------------------------------

\section[收敛关系]{收敛之间的关系}

%------------------------------------------------

\begin{frame}{强 $\Longrightarrow$ 弱}

\begin{prop}
若 $f_n \xrightarrow{\ \text{强}\ } f$ (即 $\lVert f_n - f \rVert_p \to 0$),
则 $f_n \xrightarrow{\ \text{弱}\ } f$.
\end{prop}

\startproof[bg=yellow, fg=red](证明)
由 Hölder 不等式, 对任意 $g \in L^q$, 有
\[
\Big| \int_E f_n g ~ \mathrm{d} m - \int_E f g ~ \mathrm{d} m \Big|
\leqslant \int_E |(f_n - f) g| ~ \mathrm{d} m
\leqslant \lVert f_n - f \rVert_p \cdot \lVert g \rVert_q
\longrightarrow 0 ,
\]
即 $\int_E f_n g ~ \mathrm{d} m \to \int_E f g ~ \mathrm{d} m$.
\qed

\pause

\begin{itemize}
\item \textbf{反之不成立}, 见下一页的 $\sin nx$ 例子
  (引用第 4 章习题 24, Riemann--Lebesgue 引理).
\end{itemize}

\end{frame}

%------------------------------------------------

\begin{frame}{反例: $\sin nx$ 弱收敛而不强收敛}

\textbf{例}: 由第 4 章习题 24 (\alert{Riemann--Lebesgue 引理}):
$f \in L^1(a,b)$ 时,
\[
\lim_{t \to +\infty} \int_{[a,b]} f(x)\, e^{itx} ~ \mathrm{d} x
= \lim_{t \to +\infty} \Big( \int_{[a,b]} f(x) \cos tx ~ \mathrm{d} x
+ i \int_{[a,b]} f(x) \sin tx ~ \mathrm{d} x \Big) = 0 .
\]
\pause
取 $t = n$, $n \in \mathbb{N}$, 可得到
$\lim\limits_{n \to \infty} \int_{[a,b]} f(x) \sin nx ~ \mathrm{d} x = 0$.
于是对 $g \in L^2(0, \pi) \subset L^1(0, \pi)$ ($mE = \pi < \infty$):
\[
\int_0^{\pi} g \sin nx ~ \mathrm{d} x \longrightarrow 0
\qquad \Longrightarrow \qquad
\{ \sin nx \} \xrightarrow[\ L^2(0, \pi)\ ]{\ \text{弱}\ } 0 .
\]
\pause
但 $\{ \sin nx \}$ \alert{不是 $L^2(0,\pi)$ 中的基本列}: $n \neq m$ 时 (正交性)
\[
\int_0^{\pi} \big( \sin nx - \sin mx \big)^2 ~ \mathrm{d} x
= \int_0^{\pi} \sin^2 nx ~ \mathrm{d} x
+ \int_0^{\pi} \sin^2 mx ~ \mathrm{d} x
= \frac{\pi}{2} + \frac{\pi}{2} = \pi,
\]
且指标 $\lVert \sin nx \rVert_2 = \sqrt{\pi / 2} \not\to 0$, 弱收敛 $\neq$ 强收敛.

\end{frame}

%------------------------------------------------

\begin{frame}{一致收敛与弱收敛: $mE = \infty$ 的反例}

\begin{itemize}
\item \textcircled{2} $mE < \infty$ 时: \alert{一致收敛}
  $\Longrightarrow$ \alert{强收敛} (有界收敛定理)
  $\Longrightarrow$ 弱收敛;
\item 但 $mE = \infty$ 时, 一致收敛 $\not\Longrightarrow$ 弱收敛. 反例如下.
\end{itemize}

\pause

\begin{eg}
取
\[
f_n = \frac{1}{n}\, \chi_{(1,\ e^n)}
\qquad \Longrightarrow \qquad
f_n \in L^p(\mathbb{R}) .
\]
那么 $\forall \varepsilon > 0$, 取 $N = [1/\varepsilon] + 1$,
则 $\forall n > N$, 以及 $\forall x \in \mathbb{R}$, 有
$|f_n(x) - 0| \leqslant 1/n < \varepsilon$,
故 $f_n$ \alert{一致收敛}到 $f = 0$.

\pause
但取 $g = \chi_{(1, +\infty)}(x) / x$, 那么 $g \in L^q(\mathbb{R})$, 而
\[
\int_{\mathbb{R}} f_n g ~ \mathrm{d} m
= \frac{1}{n} \int_1^{e^n} \frac{\mathrm{d} x}{x} = 1 \neq 0 ,
\]
故 $f_n$ \alert{不弱收敛}于 $0$.
\end{eg}

\end{frame}

%------------------------------------------------

\section[小结]{小结}

%------------------------------------------------

\begin{frame}{小结: 本节结论链}

\begin{itemize}
\item \textbf{可分性}: $1 \leqslant p < \infty$ 时 $L^p$ 可分:
  三步逼近 (截断 $\to$ 简单 $\to$ 有理阶梯) $+\ \varepsilon/3$ 拼接,
  稠密可列集 $S = \big\{ \sum_i r_i \chi_{[\alpha_i, \beta_i]} \big\}$;
  误差估计的功夫都在分母里: $(2M)^{p-1}$, $2n(2M)^p$, $(2mE)^{1/p}$;
\item \textbf{另一条路 (超纲)}: 正算子 + Korovkin ($1, x, x^2$ 试验函数,
  $(t-x)^2$ 技巧) $\Rightarrow$ Bernstein 定理 $\Rightarrow$
  Weierstra\ss{} 逼近 $\Rightarrow$ $L^p$ 可分 (区间情形);
\item \textbf{弱收敛}: $f_n \xrightarrow{\ \text{弱}\ } f
  \iff \forall g \in L^q,\ \int_E f_n g ~ \mathrm{d} m \to \int_E f g ~ \mathrm{d} m$;
  双线性配对给出 $L^q \subset (L^p)^*$;
\end{itemize}

\pause

\[
\text{一致收敛}
\ \xrightarrow{\ mE < \infty\ }\
\text{强收敛}
\ \xrightarrow{\ \text{Hölder}\ }\
\text{弱收敛},
\qquad \text{逆向均不成立:}
\]
\begin{itemize}
\item $\sin nx$: 弱 $\not\Rightarrow$ 强
  (非基本列, $\int_0^{\pi} (\sin nx - \sin mx)^2 ~ \mathrm{d} x = \pi$);
\item $f_n = \frac{1}{n} \chi_{(1, e^n)}$: $mE = \infty$ 时一致 $\not\Rightarrow$ 弱.
\end{itemize}

\end{frame}

%------------------------------------------------

\end{document}
